\subsubsection{Testavimo duomenų klasifikavimo rezultatai}

Žemiau pateikiama lentelė (žr.~\ref{tab:test_predictions} lentelę) su kiekvieno testavimo duomenų įrašo neurono nustatyta klase, tikrąja klase ir neapdorota sigmoidinės funkcijos~\eqref{eq:sigmoid} išėjimo reikšme. Iš $137$ testavimo įrašų $131$ buvo klasifikuotas teisingai, $6$ – neteisingai. Lentelės formatavimui buvo pasinaudota „Claude“~\cite{claude_bendrai}.

\rowcolors{1}{white}{gray!10}
\begin{longtable}{ccccc}
    \caption{Testavimo duomenų klasifikavimo rezultatai (geriausias modelis).}
    \label{tab:test_predictions} \\
    \toprule
    Nr. & Numatyta klasė & Tikra klasė & Išėjimo reikšmė & Ar teisingai? \\
    \midrule
    \endfirsthead
    
    \toprule
    Nr. & Numatyta klasė & Tikra klasė & Išėjimo reikšmė & Ar teisingai? \\
    \midrule
    \endhead
    
    \bottomrule
    \endfoot
    
    \bottomrule
    \endlastfoot

    $1$ & $0$ & $0$ & $0{,}0023$ & Taip \\
    $2$ & $0$ & $0$ & $0{,}0009$ & Taip \\
    $3$ & $0$ & $0$ & $0{,}0053$ & Taip \\
    $4$ & $1$ & $1$ & $0{,}9999$ & Taip \\
    $5$ & $0$ & $0$ & $0{,}0009$ & Taip \\
    $6$ & $0$ & $0$ & $0{,}0089$ & Taip \\
    $7$ & $0$ & $0$ & $0{,}0098$ & Taip \\
    $8$ & $1$ & $1$ & $1{,}0000$ & Taip \\
    $9$ & $0$ & $0$ & $0{,}1606$ & Taip \\
    $10$ & $0$ & $0$ & $0{,}0011$ & Taip \\
    $11$ & $0$ & $0$ & $0{,}0143$ & Taip \\
    $12$ & $0$ & $0$ & $0{,}0042$ & Taip \\
    $13$ & $1$ & $0$ & $0{,}6281$ & \textbf{Ne} \\
    $14$ & $0$ & $0$ & $0{,}0067$ & Taip \\
    $15$ & $1$ & $1$ & $0{,}9890$ & Taip \\
    $16$ & $0$ & $0$ & $0{,}0012$ & Taip \\
    $17$ & $1$ & $1$ & $0{,}9987$ & Taip \\
    $18$ & $1$ & $1$ & $0{,}9999$ & Taip \\
    $19$ & $0$ & $0$ & $0{,}0006$ & Taip \\
    $20$ & $0$ & $0$ & $0{,}0072$ & Taip \\
    $21$ & $1$ & $1$ & $1{,}0000$ & Taip \\
    $22$ & $1$ & $1$ & $0{,}9912$ & Taip \\
    $23$ & $1$ & $1$ & $0{,}9986$ & Taip \\
    $24$ & $0$ & $0$ & $0{,}0061$ & Taip \\
    $25$ & $0$ & $0$ & $0{,}0004$ & Taip \\
    $26$ & $0$ & $0$ & $0{,}0067$ & Taip \\
    $27$ & $1$ & $1$ & $0{,}9995$ & Taip \\
    $28$ & $1$ & $1$ & $0{,}9995$ & Taip \\
    $29$ & $1$ & $1$ & $1{,}0000$ & Taip \\
    $30$ & $0$ & $0$ & $0{,}0004$ & Taip \\
    $31$ & $1$ & $1$ & $0{,}9999$ & Taip \\
    $32$ & $0$ & $0$ & $0{,}0014$ & Taip \\
    $33$ & $1$ & $1$ & $1{,}0000$ & Taip \\
    $34$ & $0$ & $0$ & $0{,}0029$ & Taip \\
    $35$ & $0$ & $0$ & $0{,}0008$ & Taip \\
    $36$ & $0$ & $0$ & $0{,}0025$ & Taip \\
    $37$ & $1$ & $1$ & $1{,}0000$ & Taip \\
    $38$ & $0$ & $0$ & $0{,}0004$ & Taip \\
    $39$ & $0$ & $0$ & $0{,}0016$ & Taip \\
    $40$ & $1$ & $1$ & $0{,}9999$ & Taip \\
    $41$ & $0$ & $0$ & $0{,}0052$ & Taip \\
    $42$ & $0$ & $0$ & $0{,}0005$ & Taip \\
    $43$ & $0$ & $0$ & $0{,}0014$ & Taip \\
    $44$ & $0$ & $0$ & $0{,}0007$ & Taip \\
    $45$ & $0$ & $0$ & $0{,}0005$ & Taip \\
    $46$ & $0$ & $0$ & $0{,}0014$ & Taip \\
    $47$ & $0$ & $0$ & $0{,}0113$ & Taip \\
    $48$ & $0$ & $0$ & $0{,}0004$ & Taip \\
    $49$ & $0$ & $0$ & $0{,}0016$ & Taip \\
    $50$ & $0$ & $0$ & $0{,}0036$ & Taip \\
    $51$ & $1$ & $1$ & $1{,}0000$ & Taip \\
    $52$ & $0$ & $0$ & $0{,}0053$ & Taip \\
    $53$ & $1$ & $1$ & $1{,}0000$ & Taip \\
    $54$ & $0$ & $0$ & $0{,}0017$ & Taip \\
    $55$ & $0$ & $0$ & $0{,}0003$ & Taip \\
    $56$ & $0$ & $0$ & $0{,}0006$ & Taip \\
    $57$ & $0$ & $1$ & $0{,}0343$ & \textbf{Ne} \\
    $58$ & $1$ & $1$ & $0{,}9997$ & Taip \\
    $59$ & $0$ & $0$ & $0{,}0011$ & Taip \\
    $60$ & $0$ & $0$ & $0{,}0006$ & Taip \\
    $61$ & $1$ & $0$ & $0{,}9960$ & \textbf{Ne} \\
    $62$ & $0$ & $0$ & $0{,}0033$ & Taip \\
    $63$ & $0$ & $0$ & $0{,}0014$ & Taip \\
    $64$ & $1$ & $1$ & $1{,}0000$ & Taip \\
    $65$ & $0$ & $0$ & $0{,}0006$ & Taip \\
    $66$ & $0$ & $0$ & $0{,}0005$ & Taip \\
    $67$ & $1$ & $1$ & $1{,}0000$ & Taip \\
    $68$ & $1$ & $1$ & $1{,}0000$ & Taip \\
    $69$ & $1$ & $1$ & $1{,}0000$ & Taip \\
    $70$ & $1$ & $1$ & $0{,}9997$ & Taip \\
    $71$ & $1$ & $1$ & $0{,}9997$ & Taip \\
    $72$ & $0$ & $0$ & $0{,}0078$ & Taip \\
    $73$ & $0$ & $0$ & $0{,}0058$ & Taip \\
    $74$ & $1$ & $0$ & $0{,}9859$ & \textbf{Ne} \\
    $75$ & $0$ & $0$ & $0{,}0029$ & Taip \\
    $76$ & $0$ & $0$ & $0{,}0053$ & Taip \\
    $77$ & $0$ & $0$ & $0{,}0114$ & Taip \\
    $78$ & $1$ & $1$ & $1{,}0000$ & Taip \\
    $79$ & $0$ & $0$ & $0{,}0008$ & Taip \\
    $80$ & $1$ & $1$ & $0{,}9998$ & Taip \\
    $81$ & $0$ & $0$ & $0{,}0017$ & Taip \\
    $82$ & $0$ & $0$ & $0{,}0007$ & Taip \\
    $83$ & $0$ & $0$ & $0{,}0004$ & Taip \\
    $84$ & $0$ & $0$ & $0{,}0005$ & Taip \\
    $85$ & $1$ & $1$ & $1{,}0000$ & Taip \\
    $86$ & $0$ & $0$ & $0{,}0006$ & Taip \\
    $87$ & $1$ & $1$ & $0{,}9995$ & Taip \\
    $88$ & $0$ & $0$ & $0{,}0237$ & Taip \\
    $89$ & $0$ & $0$ & $0{,}0016$ & Taip \\
    $90$ & $1$ & $1$ & $0{,}9997$ & Taip \\
    $91$ & $0$ & $0$ & $0{,}0062$ & Taip \\
    $92$ & $0$ & $0$ & $0{,}0012$ & Taip \\
    $93$ & $1$ & $1$ & $0{,}9883$ & Taip \\
    $94$ & $0$ & $0$ & $0{,}0005$ & Taip \\
    $95$ & $0$ & $0$ & $0{,}0033$ & Taip \\
    $96$ & $0$ & $0$ & $0{,}0029$ & Taip \\
    $97$ & $0$ & $0$ & $0{,}0098$ & Taip \\
    $98$ & $1$ & $1$ & $0{,}9991$ & Taip \\
    $99$ & $1$ & $1$ & $1{,}0000$ & Taip \\
    $100$ & $1$ & $1$ & $1{,}0000$ & Taip \\
    $101$ & $0$ & $0$ & $0{,}0005$ & Taip \\
    $102$ & $0$ & $0$ & $0{,}0036$ & Taip \\
    $103$ & $0$ & $0$ & $0{,}0004$ & Taip \\
    $104$ & $0$ & $0$ & $0{,}0016$ & Taip \\
    $105$ & $1$ & $1$ & $0{,}9677$ & Taip \\
    $106$ & $0$ & $1$ & $0{,}2578$ & \textbf{Ne} \\
    $107$ & $0$ & $0$ & $0{,}0098$ & Taip \\
    $108$ & $1$ & $1$ & $1{,}0000$ & Taip \\
    $109$ & $1$ & $1$ & $1{,}0000$ & Taip \\
    $110$ & $0$ & $0$ & $0{,}0005$ & Taip \\
    $111$ & $0$ & $0$ & $0{,}0042$ & Taip \\
    $112$ & $1$ & $1$ & $1{,}0000$ & Taip \\
    $113$ & $1$ & $1$ & $1{,}0000$ & Taip \\
    $114$ & $0$ & $0$ & $0{,}0145$ & Taip \\
    $115$ & $0$ & $0$ & $0{,}0013$ & Taip \\
    $116$ & $0$ & $0$ & $0{,}0005$ & Taip \\
    $117$ & $1$ & $1$ & $1{,}0000$ & Taip \\
    $118$ & $1$ & $1$ & $1{,}0000$ & Taip \\
    $119$ & $0$ & $0$ & $0{,}0020$ & Taip \\
    $120$ & $0$ & $0$ & $0{,}0042$ & Taip \\
    $121$ & $1$ & $1$ & $1{,}0000$ & Taip \\
    $122$ & $0$ & $0$ & $0{,}0007$ & Taip \\
    $123$ & $0$ & $1$ & $0{,}3784$ & \textbf{Ne} \\
    $124$ & $1$ & $1$ & $0{,}9990$ & Taip \\
    $125$ & $0$ & $0$ & $0{,}0012$ & Taip \\
    $126$ & $0$ & $0$ & $0{,}0023$ & Taip \\
    $127$ & $1$ & $1$ & $1{,}0000$ & Taip \\
    $128$ & $0$ & $0$ & $0{,}0007$ & Taip \\
    $129$ & $0$ & $0$ & $0{,}0381$ & Taip \\
    $130$ & $1$ & $1$ & $0{,}9816$ & Taip \\
    $131$ & $0$ & $0$ & $0{,}0279$ & Taip \\
    $132$ & $1$ & $1$ & $0{,}9998$ & Taip \\
    $133$ & $1$ & $1$ & $1{,}0000$ & Taip \\
    $134$ & $0$ & $0$ & $0{,}0010$ & Taip \\
    $135$ & $1$ & $1$ & $1{,}0000$ & Taip \\
    $136$ & $0$ & $0$ & $0{,}0045$ & Taip \\
    $137$ & $1$ & $1$ & $1{,}0000$ & Taip \\
\end{longtable}

Iš $137$ testavimo įrašų $131$ buvo klasifikuotas teisingai ($95{,}6\ \%$ tikslumas). Neteisingai klasifikuoti $6$ įrašai:
\begin{itemize}
    \item $2$ nepiktybiniai navikai (klasė $0$) buvo priskirti piktybiniams (klasė $1$) – įrašai Nr. $13$ ir $74$.
    \item $2$ piktybiniai navikai (klasė $1$) buvo priskirti nepiktybiniams (klasė $0$) – įrašai Nr. $57$ ir $123$.
    \item $2$ nepiktybiniai navikai (klasė $0$) buvo priskirti piktybiniams (klasė $1$) – įrašas Nr. $61$, ir $1$ piktybinis (klasė $1$) priskirtas nepiktybiniams – įrašas Nr. $106$.
\end{itemize}